\documentclass[11pt]{article}
\usepackage{amsfonts, array,enumitem,natbib,bibentry,amsmath}
\bibliographystyle{plain}
\usepackage[T1]{fontenc}
\usepackage{charter,euler}
\usepackage[margin=0.6in]{geometry}
\usepackage[hidelinks]{hyperref}
\hypersetup{
  breaklinks=true,
  colorlinks=true,
  urlcolor=blue
}
%\setlist[itemize]{leftmargin=*}
\usepackage{titlesec}
\titleformat{\section}{\bfseries \large\scshape\raggedright}{}{0em}{}
[\titlerule]
\titlespacing\section{0pt}{0pt}{7pt}
\newif\iflong
\longtrue % set to true if cv, false if resume

\newif\ifpublic
\publictrue % set to true for the published CV (suppresses pending/declined grants and
            % any other private content). Set to false for the author's full personal copy.
%--------------------BEGIN DOCUMENT----------------------
\begin{document}
%\nobibliography{cv}

%\includepdf[pages={1-3}]{cv.pdf}
\pagestyle{empty} % non-numbered pages

%--------------------TITLE------------
\begin{center}
\textbf{\Large Dr. Sushovan Majhi}\\
Assistant Professor of Data Science, George Washington University\\
Samson Hall 313, George Washington University, Washington, D.C., USA\\
\href{mailto:s.majhi@gwu.edu}{\tt s.majhi@gwu.edu} |
\href{tel:+12029941235}{\tt (+1) 202-994-1235} |
\href{https://www.smajhi.com}{\tt Homepage} |
\href{https://www.github.com/sushovan4}{\tt GitHub} |
\href{https://scholar.google.com/citations?user=3zP2QQsAAAAJ&hl=en&oi=sra}
{\tt Google Scholar} |
\href{https://www.linkedin.com/in/smajhi/}{\tt LinkedIn}
\end{center}
%------------------------------------
\section{\textsc{Research Interests}} \textbf{Applied Topology}, 
\textbf{Topological Data Analysis (TDA)},
\textbf{Discrete and Computational Geometry}, \textbf{Shape Matching}, \textbf{Statistical Finance}, and \textbf{Topological Deep Learning}.

My research primarily revolves around the interface of
\emph{mathematics}, \emph{computer science}, and the mathematical foundations of \emph{data science}. 
More specifically, I am motivated to develop provable inference techniques for data science that are inspired by topology and geometry. I also keep a keen interest in applying TDA to fascinating, real-world problems arising in fields, like finance,  dynamical systems, biology, medicine, and genetics. \\

\section{Education}
\begin{itemize}[leftmargin=*]
\item \textbf{Doctor of Philosophy} in \textbf{Mathematics} 
  \hfill\textbf{\footnotesize August 2014--December 2020}\\
  \textbf{Tulane University}, New Orleans, LA, USA.\\
  \emph{Advisor}: \textbf{Prof. Carola Wenk} \\
  \emph{Courses}: computational geometry, computational topology,
  topological data analysis,  differential geometry, differentiable manifolds,
  algorithms, data structures, computational complexity, applied mathematics, and scientific computing.

\item  \textbf{Master of Science} in \textbf{Mathematics}
  \hfill\textbf{\footnotesize August 2009--May 2012}\\
  \textbf{Tata Institute of Fundamental Research}, Bangalore, India \\
  \emph{Courses:} ordinary and partial differential equations, probability
  theory, complex analysis, functional analysis, numerical linear algebra,
  measure theory, mechanics.
  
\item \textbf{Bachelor of Science} in \textbf{Mathematics} (Hons.)
  \hfill\textbf{\footnotesize July 2006--May 2009}\\
  \textbf{Ramakrishna Mission Vidyamandira}, \textbf{Calcutta University},
  West Bengal, India\\
  \emph{Courses}: calculus, real analysis, linear algebra, numerical analysis,
  game theory, statistics, physics.
\end{itemize}

\section{Work Experience}
\begin{itemize}[leftmargin=*]
\item \textbf{Assistant Professor} 
\hfill\textbf{\footnotesize August 2024--\emph{current}}\\
\textbf{Data Science Program, George Washington University}, Washington D.C., USA
\item \textbf{Visiting Assistant Professor} 
\hfill\textbf{\footnotesize August 2023--July 2024}\\
\textbf{Data Science Program, George Washington University}, Washington D.C., USA\\
\item \textbf{Postdoctoral Research Fellow} 
\hfill\textbf{\footnotesize January 2021--July 2023}\\
\textbf{School of Information, University of California}, Berkeley, USA\\
\textit{Role}: The responsibilities included conducting research broadly in data
science, forging new research collaborations, organizing research webinar
series.
\end{itemize}

\section{Preprints}
\begin{enumerate}
\item[9.] \textbf{SM}, Atish Mitra, \v{Z}iga Virk, Pramita Bagchi. A Closed-Form Adaptive-Landmark Kernel for Certified Point-Cloud and Graph Classification, May 2026.
Available at:
\href{https://arxiv.org/abs/2605.04046}{\tt arXiv:2605.04046 [cs.LG]}\\
$\bullet$ Submitted to \emph{Journal of Machine Learning Research}

\item[8.] \textbf{SM}, Atish Mitra, \v{Z}iga Virk, Pramita Bagchi. A Closed-Form Persistence-Landmark Pipeline for Certified Point-Cloud and Graph Classification, May 2026.
Available at:
\href{https://arxiv.org/abs/2605.02836}{\tt arXiv:2605.02836 [cs.LG]}\\
$\bullet$ Submitted to \emph{Transactions on Machine Learning Research}

\item[7.] Salam Rabindrajit Luwang, \textbf{SM}, Vishal Mandal, Atish J. Mitra, Md. Nurujjaman, Buddha Nath Sharma. Interpretable Classification of Time Series Using Euler Characteristic Surfaces
, 25 pages, March 2026.
Available at:
\href{https://arxiv.org/abs/2603.15079}{\tt 	arXiv:2603.15079 [cs.LG]}

\item[6.] Kazuhiro Kawamura, \textbf{SM}, Atish Mitra. The Shadow of Vietoris–Rips Complexes in Limits, 29 pages, January 2026.
Available at:
\href{https://arxiv.org/abs/2601.01359}{\tt arXiv:2601.01359 [math.AT]}\\
$\bullet$ Submitted to \emph{Journal of Applied and Computational Topology}

\item[5.] Rafal Komendarczyk, \textbf{SM}, Atish Mitra. Vietoris--Rips Shadow for Euclidean Graph Reconstruction, 19 pages, June 2025.
Available at:
\href{https://arxiv.org/abs/2506.01603}{\tt arXiv:2506.01603 [math.AT]}\\
$\bullet$ Submitted to \emph{Journal of Applied and Computational Topology}

\item[4.] Joshua Dorrington, \textbf{SM}, Atish Mitra, James Moukheiber*, Demi Qin, Jacob Sriraman*, Kristian Strommen. Topology of the Polar Vortex and Montana Weather, 6 pages, Mar 2025.
Available at:
\href{https://arxiv.org/abs/2503.20743}{\tt arXiv:2503.20743 [math.DS]}\\
$\bullet$ Presented at \emph{ATMCS'25} \& subitted to \emph{La Matematica}

\item[3.] Henry Adams, \textbf{SM}, Fedya Manin, Nicol\`o Zava, \v{Z}iga Virk. Lower Bounding The Gromov--Hausdorff Distance In Metric Graphs, 20 pages, Nov 2024.
Available at:
\href{https://arxiv.org/abs/2411.09182}{\tt arXiv:2411.09182v1 [math.MG]}\\
$\bullet$ Accepted to SoCG'26 \& submitted to \emph{Discrete and Computational Geometry}


%Available at: \href{https://arxiv.org/abs/2406.04259}{\tt arXiv:2406.04259 [math.AT]}
\item [2.] Halley Fritze, \textbf{SM}, Marissa Masden, Atish Mitra, Michael Stickney. Faithful Reeb Graph Reconstruction of a Tectonic Subduction Zone from Earthquake Hypocenters. 
Available at:
\href{https://arxiv.org/abs/2410.19410}{\tt arXiv:2410.19410 [cs.CG]}\\
$\bullet$ Presented at \emph{ATMCS'25} \& submitted to \emph{La Matematica}


\item[1.] Rafal Komendarczyk, \textbf{SM}, Will Tran. Topological Stability and Latschev-type Reconstruction Theorems for Spaces of Curvature Bounded Above, 27 pages, June 2024. Available at:
\href{https://arxiv.org/abs/2406.04259}{\tt arXiv:2406.04259 [math.AT]}\\
$\bullet$ Submitted to \emph{Discrete \& Computational Geometry}

\end{enumerate}

\let\thefootnote\relax\footnote{* Undergraduate student at the time of research or submission}

\section{Refereed Journal Publications}
\begin{enumerate}
\item[10.] Buddha Nath Sharma, Anish Rai, SR Luwang, Md. Nurujjaman, \textbf{SM}. Causality Analysis of COVID-19 Induced Crashes in Stock and Commodity Markets: A Topological Perspective. \emph{International Journal of Modern Physics C}, September 2025.
DOI: \href{https://doi.org/10.1142/S0129183126500221}{\tt 10.1142/S0129183126500221}.
Also available at:
\href{https://doi.org/10.48550/arXiv.2502.14431}{\tt arXiv:2502.14431 [q-fin.ST]}

\item[9.] Henry Adams, Florian Frick, \textbf{SM}, Nichola McBride\textsuperscript{*}. Hausdorff VS Gromov--Hausdorff Distances.
\emph{Discrete and Computational Geometry}, February 2025.
DOI: \href{https://doi.org/10.1007/s00454-025-00722-9}{\tt 10.1007/s00454-025-00722-9}
Also available at:
\href{https://arxiv.org/abs/2309.16648}{\tt arXiv:2309.16648[math.MG]}

\item[8.] Kundan Mukhia, Anish Rai, SR Luwang, Md Nurujjaman, \textbf{SM}, Chittaranjan Hens. Complex Network Analysis of Cryptocurrency Market During Crashes. \emph{Physica A: Statistical Mechanics and its Applications}
, November 2024. 
DOI: \href{https://doi.org/10.1016/j.physa.2024.130095}{\tt 10.1016/j.physa.2024.130095}.
Also available at:
\href{https://arxiv.org/abs/2405.05642}{\tt arXiv:2405.05642 [q-fin.ST]}

\item[7.] Anish Rai, Buddha Nath Sharma, Salam Rabindrajit Luwang, Md.Nurujjaman, \textbf{SM}. Identifying Extreme Events in the Stock Market: A Topological Data Analysis. \emph{Chaos: An Interdisciplinary Journal of Nonlinear Science},
October 2024. 
DOI: \href{https://doi.org/10.1063/5.0220424}{\tt 10.1063/5.0220424}. 
Also available at:
\href{https://arxiv.org/abs/2405.16052}{\tt arXiv:2405.16052 [q-fin.ST]}\\


\item[6.] \textbf{SM}. Demystifying Latschev's theorem: Manifold Reconstruction
from noisy data. \emph{Discrete \& Computational Geometry}, May 2024. 
DOI: \href{https://doi.org/10.1007/s00454-024-00655-9}{\tt 10.1007/s00454-024-00655-9}. 
Also available at:
\href{https://arxiv.org/abs/2305.17288}{\tt arXiv:2305.17288[math.AT]}

\item[5.] \textbf{SM} and Carola Wenk. Distance Measures for geometric graphs.
\emph{Computational Geometry: Theory and Applications}, March 2024. 
DOI: \href{https://doi.org/10.1016/j.comgeo.2023.102056}{\tt 10.1016/j.comgeo.2023.102056}. Also available at: \href{https://arxiv.org/abs/2209.12869}{\tt arXiv:2209.12869[cs.CG]}

\item[4.] \textbf{SM}. Vietoris--Rips complexes of metric spaces near a 
metric graph. \emph{Journal of Applied and Computational Topology}, May 2023.
DOI: \href{https://doi.org/10.1007/s41468-023-00122-z}{\tt 10.1007/s41468-023-00122-z}.
Available at: \href{https://arxiv.org/abs/2204.14234}{\tt arXiv:2204.14234[math.AT]}

\item[3.] \textbf{SM}, Jeffrey Vitter, and Carola Wenk. 
Approximating Gromov-Hausdorff distance in Euclidean space. 
\emph{Computational Geometry: Theory and Applications}, January 2024.
DOI: \href{https://doi.org/10.1016/j.comgeo.2023.102034}{\tt 10.1016/j.comgeo.2023.102034}.
Also available at: \href{https://arxiv.org/abs/1912.13008}{\tt arXiv:1912.13008[math.MG]}

\item[2.] Anish Rai, Ajit Mahata, Md Nurujjaman, \textbf{SM}, and Kanish Debnath. A sentiment-based 
modeling and analysis of stock price during the COVID-19: U- and Swoosh-shaped recovery. 
\emph{Physica A: Statistical Mechanics and its Applications}. April 2022.
DOI: \href{https://www.sciencedirect.com/science/article/pii/S0378437121009778}{\tt 10.1016/j.physa.2021.126810}. Also available at: \href{https://arxiv.org/abs/2110.03986}{\tt arXiv:2110.03986 [q-fin.ST]}

\item[1.] Brittany Terese Fasy, Rafal Komendarczyk, \textbf{SM}, and Carola Wenk. On the reconstruction of
geodesic subspaces of $\mathbb R^N$. \emph{International Journal of Computational Geometry \& Applications}. September
2022. DOI: \href{https://www.worldscientific.com/doi/10.1142/S0218195922500066}{\tt 10.1142/S0218195922500066}.
Also available at: \href{https://arxiv.org/abs/1810.10144}{\tt arXiv:1810.10144[math.AT]}
\end{enumerate}

\section{Refereed Conference Proceedings}
\begin{itemize}
\item[4.] Henry Adams, \textbf{SM}, Fedya Manin, Nicol\`o Zava, \v{Z}iga Virk. Lower Bounding The Gromov--Hausdorff Distance In Metric Graphs.
In Proceedings of the \emph{42nd International Symposium on Computational Geometry (SoCG)}, New Jersey, USA, June 2–5, 2026.
Available at:
\href{https://arxiv.org/abs/2411.09182}{\tt arXiv:2411.09182v1 [math.MG]}
%\\
%$\bullet$ Submitted to \emph{Discrete and Computational Geometry}

\item[3.] \textbf{SM}. Demystifying Latschev's theorem: Manifold Reconstruction
from noisy data. In Proceedings of
the \emph{40th International Symposium on Computational Geometry (SoCG)}, Athens, Greece, June 11–14, 2024. \\ 
Available at:
\href{https://doi.org/10.4230/LIPIcs.SoCG.2024.73}{\tt https://doi.org/10.4230/LIPIcs.SoCG.2024.73}

\item[2.] Erin Chambers, Brittany Fasy, Benjamin Holmgren\textsuperscript{*},
\textbf{SM}, and Carola Wenk. Metric and path-connectedness
properties of the Fr\v{e}chet distance for paths and graphs. In Proceedings of
the \emph{34th Canadian Conference on Computational Geometry (CCCG)}, 2023. DBLP: \href{https://dblp.org/rec/conf/cccg/ChambersHFMW23.html}{\tt conf/cccg/ChambersHFMW23}.
Available at:
\href{https://arxiv.org/abs/2308.00900}{\tt arXiv:2308.00900[cs.CG]}
  
\item[1.] \textbf{SM}. Graph mover's distance: An efficiently
computable distance measure for geometric graphs. In Proceedings of the \emph{34th
Canadian Conference on Computational Geometry (CCCG)}, 2023.\\
DBLP: \href{https://dblp.org/rec/conf/cccg/Majhi23.html}{\tt conf/cccg/Majhi23}. 
Available at: \href{https://arxiv.org/abs/2306.02133}{\tt arXiv:2306.02133[cs.CG]}
\end{itemize}


\section{Workshop Contributions}
\begin{itemize}
\item[8.] Pramita Bagchi, \textbf{SM}, Atish Mitra, \v{Z}iga Virk. Statistical Analysis of Persistence Landmarks: A Hilbert Space
Embedding of the Space of Persistence Diagrams, At \emph{Fall Workshop on Computational Geometry (FWCG)}, Queens College, 2025. 

\item[7.] Enrique Alvarado, Daniela Beckelhymer, Joshua Dorrington, Tung Lam, \textbf{SM}, Jasmine Noory, María Sánchez Muniz, Kristian Strommen. Detecting the Indian Monsoon using Topological Data Analysis. At \emph{ATMCS'25}. June 2025.
Available at:
\href{https://arxiv.org/abs/2504.01022}{\tt arXiv:2504.01022 [physics.ao-ph]}

\item[6.] Halley Fritze, \textbf{SM}, Marissa Masden, Atish Mitra. Hausdorff--Approximation of Metric Graphs, At \emph{Fall Workshop on Computational Geometry (FWCG)}, Tufts University, 2024. 
Available at:
\href{https://arxiv.org/abs/2410.19410}{\tt arXiv:2410.19410 [cs.CG]}

\item[5.] Halley Fritze, \textbf{SM}, Marissa Masden, Atish Mitra. Hausdorff--Approximation of Metric Graphs, At \emph{Fall Workshop on Computational Geometry (FWCG)}, Tufts University, 2024. 
Available at:
\href{https://arxiv.org/abs/2410.19410}{\tt arXiv:2410.19410 [cs.CG]}

\item[4.] E. Chambers, B. Fasy, B. Holmgren* , \textbf{SM}, and C. Wenk.
Path-Connectivity of Fr\'echet Spaces of Graphs. Computational Geometry: Young
Researchers Forum, 2022

\item[3.] \textbf{SM} and Carola Wenk. Distance Measures for Geometric
Graphs, At \emph{Fall Workshop on Computational Geometry (FWCG)}, 2022
  
\item[2.] Brittany Terese Fasy, \textbf{SM} and Carola Wenk. Threshold-based
  graph reconstruction using discrete Morse theory. In \emph{Fall Workshop on
  Computational Geometry (FWCG)}, 2018. Abstract available at:
  \href{https://arxiv.org/abs/1911.12865}{\tt Link}
 
\item[1.] Brittany Terese Fasy, Rafal Komendaczyk, \textbf{SM}, and Carola
  Wenk. Topological reconstruction of metric graphs in $\mathbb R^N$. At
  \emph{Fall Workshop on Computational Geometry (FWCG)}, 2017. Abstract
  available at:
  \href{http://www.ams.stonybrook.edu/~jsbm/fwcg17/proceedings.pdf} {\tt Link}
\end{itemize}

%\let\thefootnote\relax\footnote{* Undergraduate student at the time of research or submission}

\section{Graduate Teaching Experience}
\begin{itemize}[leftmargin=*]

\item \textbf{Linear Algebra for Data Science} \hfill
  \textbf{\footnotesize Summer 2024}\\
  Data Science Program, George Washington University, USA \\
  \emph{Topics:} linear algebra foundations, Gaussian elimination, pivoting techniques for stability, LU decomposition, eigenvalues and eigenvectors,
  solving linear systems using iterative methods, projections and linear regression, SVD and applications.
  \emph{Website:} \href{https://linalg.mathematics.land/}{\tt linalg.mathematics.land}

\item \textbf{Algorithm Design (DATS 6001)} \hfill
  \textbf{\footnotesize Spring 2024--\emph{current}}\\
  Data Science Program, George Washington University, USA \\
  \emph{Topics:} basics of computational complexity, data structures (array,
  stack, queue, tree, etc), algorithms (search, sort, etc), and programming
  paradigms like dynamic programming and greedy algorithms.\\
    \emph{Website:} \href{https://6001-majhi.datasci.land/}{\tt 6001-majhi.datasci.land}

\item \textbf{Topological Data Analysis (CS)} \hfill
  \textbf{\footnotesize Fall 2023}\\
  National Institute of Technology, Sikkim, India \\
  \emph{Topics:} topological spaces, metric spaces and
  their examples, simplicial complexes, homology in $\mathbb Z/2$ coeff,
  persistent homology, bottleneck and Wasserstein distance, and non-linear
  time-series analysis using TDA. \emph{Website:} \href{https://tda.mathematics.land/}{\tt tda.mathematics.land}

\item \textbf{Computer Science Foundations (DATS 6450)} \hfill
  \textbf{\footnotesize Fall 2023}\\
  Data Science Program, George Washington University, USA \\
  \emph{Topics:} computer design, programming in Python, object-oriented programming

\item \textbf{Introduction to Data Mining (DATS 6103)} \hfill
  \textbf{\footnotesize Fall 2023--\emph{current}}\\
  Data Science Program, George Washington University, USA \\
  \emph{Topics:} data wrangling, linear and logistic
  regression, classification, clustering, data visualization, support vector
  machines, machine learning algorithms.

\item \textbf{Statistics for Data Science (MIDS 203)} \hfill
  \textbf{\footnotesize August 2020--July 2023}\\
  School of Information, University of California, Berkeley, USA \\
  \emph{Topics:} probability theory,
  sampling distributions, estimators and convergence theorems, 
  confidence intervals, hypothesis testing, and regression.

\item \textbf{Linear Algebra, Complex Analysis}
  \hfill \textbf{\footnotesize January 2013--June 2013}\\
  Christ University and Scimetric Pvt Ltd, Bangalore, India
 
\item  \textbf{Analysis, Linear Algebra, Complex Analysis}
  \hfill  \textbf{\footnotesize November 2011--July 2012} \\
  GATE-IIT Coaching Institute, JP Nagar, Bangalore, India \\
  Graduate level,  for competitive national exams, e.g., National Eligibility Test 
\item \textbf{Analysis and Linear Algebra}
  \hfill \textbf{\footnotesize February 2012--July 2012} \\
  MES College, Department of Mathematics, Malleswaram, Bangalore, India 
\end{itemize}

\section{Undergraduate Teaching Experience}
\begin{itemize}[leftmargin=*]
\item \textbf{Data Science for All} \hfill
  \textbf{\footnotesize Summer 2025}\\
  Data Science Program, George Washington University, USA \\
  \emph{Topics:} Foundations of data science for undergraduate students

\item \textbf{Data Science Capstone (DATS 4001)} \hfill
  \textbf{\footnotesize Fall 2024--\emph{current}}\\
    Data Science Program, George Washington University, USA\\
  \emph{Goal:} The goal of the capstone is for the students to apply the knowledge acquired throughout the GWU Data Science undergraduate program to solve a real-world problem, using real data. It typically starts with developing S.M.A.R.T. questions and goes through the entire data science life cycle in multiple iterations.


\item \textbf{Undergraduate Statistics for Business Students (MATH 1140)} \hfill
  \textbf{\footnotesize Summer 2019}\\
  Tulane University, USA \\
  \emph{Topics:} sampling methods, descriptive statistics, probability theory,
  random variables, limit theorems, confidence intervals, hypothesis
  testing, and linear regression.
\end{itemize}

\section{Other Teaching}
\begin{itemize}[leftmargin=*]
\item \textbf{Graduate Teaching Assistant} \hfill
  \textbf{\footnotesize Fall 2014--Spring 2017}\\
  Tulane University, USA
\end{itemize}

\section{Invited Talks and Accepted Abstracts}
\begin{itemize}[leftmargin=*]
  \item \textbf{Title:} Gemetric Latschev's Theorem\hfill  \textbf{\footnotesize March 11, 2026}\\
  Spring Topology Conference, University of Alabama, D.C.   

  \item \textbf{Title:} Vietoris--Rips Shadow for Euclidean Graph Reconstruction\hfill  \textbf{\footnotesize January 7, 2026}\\
  Joint Mathematics Meetings, Washington, D.C.   

  \item \textbf{Title:} Vietoris--Rips Shadow for Euclidean Graph Reconstruction\hfill  \textbf{\footnotesize Nov 8, 2025}\\
  Fall Workshop on Computational Geometry, Queens College   
  
  
  \item \textbf{Title:} Vietoris--Rips Shadow for Euclidean Graph Reconstruction\hfill  \textbf{\footnotesize Oct 4, 2025}\\
  AMS Southeastern Sectional Meeting, Tulane University 

  \item \textbf{Title:} Topological and Geometric Reconstruction of Spaces Beyond Manifolds\hfill  \textbf{\footnotesize June 19, 2025}\\
  TDA Week, Kyoto University, Japan, 2025 

  \item \textbf{Title:} Topological Stability and Latschev-type Reconstruction Theorems for CAT($\kappa$)
 Spaces\hfill  \textbf{\footnotesize March 8, 2025}\\
  AMS Southeastern Sectional Meeting, Clemson University 

  \item \textbf{Title:} Topological Stability and Latschev-type Reconstruction Theorems for Spaces of Curvature Bounded Above\hfill  \textbf{\footnotesize March 6, 2025}\\
  Spring Topology and Dynamics, Christopher Newport University, Newport News 

  \item \textbf{Title:} Predicting the Onset and Withdrawal of the Indian Monsoon\\ using Persistent Homology\hfill  \textbf{\footnotesize January 9, 2025}\\
  AMS Special Session: Climate Science at the Interface Between Topological Data Analysis and Dynamical Systems Theory,
  Joint Mathematics Meetings, Seattle, 2025 

  \item \textbf{Title:} Lower Bounding the Gromov--Hausdorff distance in Metric Graphs\hfill  \textbf{\footnotesize Nov 14--15, 2024}\\
  The 31st Fall Workshop on Computational Geometry, Tufts University

  \item \textbf{Title:} A Taste of Topological Data Analysis (TDA): Reconstruction of Shapes\hfill  \textbf{\footnotesize Oct 14, 2024}\\
  Indian Institute of Technology, Mandi, India

  \item \textbf{Title:} A Taste of Topological Data Analysis (TDA): Reconstruction of Shapes\hfill  \textbf{\footnotesize Oct 8, 2024}\\
  Vellore Institute of Technology, Chennai, India

  \item \textbf{Title:} Demystifying Latschev's Theorem\hfill  \textbf{\footnotesize June 14, 2024}\\
  Symposium on Computational Geometry, Athens, Greece
  
  \item \textbf{Title:} Demystifying Latschev's Theorem\hfill  \textbf{\footnotesize March 28, 2024}\\
  Montana State University

  \item \textbf{Title:} Demystifying Latschev's Theorem\hfill  \textbf{\footnotesize August 23, 2023}\\
  Applied Algebraic Topology Research Network (AATRN)\\
  Available on: \href{https://www.youtube.com/watch?v=bSFkxlYc4P0}{\tt YouTube}

  \item \textbf{Title:} Graph Mover's Distance
  \hfill  \textbf{\footnotesize August 2--4, 2023} \\
  The 34th Canadian Conference on Computational Geometry, Montreal, Canada
 
  \item \textbf{Title:} Similarity Measures for Geometric Graphs
  \hfill  \textbf{\footnotesize October 14--15, 2022} \\
  The 30th Fall Workshop on Computational Geometry\\North Carolina State University, USA

\item \textbf{Title:} Topological Methods in the Reconstruction and Comparison
of Shapes \hfill  \textbf{\footnotesize February, 2022} \\
  Mathematics Department, ICFAI University, Tripura, India

  \item \textbf{Title:} A Taste of Topological Data Analysis (TDA): Reconstruction of Shapes
  \hfill  \textbf{\footnotesize September, 2021} \\
  Department of Mathematics, Hunter College, NY, USA

  \item \textbf{Fall Workshop on Computational Geometry}
  \hfill \textbf{\footnotesize October, 2018} \\
  Queens College, New York, USA

\item \textbf{Fall Workshop on Computational Geometry}
  \hfill \textbf{\footnotesize November, 2017}\\
  SUNY (Stony Brook), New York, USA
\end{itemize}

\section{PhD Dissertation Committee}
\begin{itemize}
  \item Peiqi Yang, Department of Mathematics, George Washington University \hfill\textbf{\footnotesize Summer 2025}

  \item Will Tran, Department of Mathematics, Tulane University, \hfill\textbf{\footnotesize Spring 2024}\\
  \emph{Thesis Title:} Distortion and Curvature in the Shape Reconstruction Problem
\end{itemize}

\section{PhD Advisees}
\begin{enumerate}
\item Buddha Nath Sharma, PhD student at \hfill\textbf{\footnotesize August 2025--\emph{current}}\\ 
National Institute of Technology, Sikkim, India \\
\emph{Thesis Title:} Applications of topological data analysis in financial markets\\
\emph{Co-advised} with \href{https://physics.nitsikkim.ac.in/profile.php?sn=23}{Md. Nurujjaman}
\end{enumerate}

\section{Graduate Student Mentees}
\begin{itemize}
\item Mauricio Montes, current PhD student\\ Auburn University\hfill\textbf{\footnotesize March 2025--August 2025}\\
\emph{Project Title:} Statistical Analysis of Euler Characteristic Surfaces

\item Khush Shah and Shikha Kumari, was data science master's students\\ George Washington University, \hfill\textbf{\footnotesize August 2023--December 2024}\\
\emph{Project Title:} Matching geometric graphs using Graph Neural Networks\\
\emph{Now:} Junior Data Scientist at Cintra and AI/ML \& Software Engineer at Nexas 8 International, respectively

\item Anish Rai, was a PhD student at\hfill\textbf{\footnotesize August 2022--August 2024}\\ 
National Institute of Technology, Sikkim, India \\
\emph{Project Title:}  Prediction of stock market crashes using topological data analysis \\
\emph{Now:} Postdoctoral scholar at Chennai Mathematical Institute (CMI), Chennai, Bangalore 
\end{itemize}

\section{Undergraduate Student Mentees}
\begin{itemize}
\item Abby Stein\\ Senior at George Washington University \hfill\textbf{\footnotesize August 2025--\emph{current}} %\\
%\emph{Now:} Master's student at Montana State University

\item Jacob Sriraman\\ was at Montana State University \hfill\textbf{\footnotesize August 2025--\emph{current}}\\
\emph{Now:} Master's student at Montana State University

\item James Moukheiber\\ was at George Washington University \hfill\textbf{\footnotesize August 2024--\emph{current}}\\
\emph{Now:} Master's student at the University of Zurich

\end{itemize}

\section{Department/University Service}
\begin{itemize}

  \item Chair, Undergraduate Curriculum Committee\\ George Washington University, Data Science Program \hfill\textbf{\footnotesize August 2025--\emph{current}}
  
  \item Undergraduate Advisor\\ George Washington University, Data Science Program \hfill\textbf{\footnotesize August 2024--\emph{current}}
\end{itemize}


\section{Awards, Scholarships, Recognition}
\begin{itemize}[leftmargin=*]
\ifpublic\else
\item NSF: (DMS) Computational Mathematics, \textbf{Pending} \\
\emph{Title}: Shape Learning from Euclidean Shadows: A Topological, Geometric, and Statistical Framework\\
\emph{Role}: PI

\item NSF: (CCF) Algorithmic Foundations, \textbf{Pending} \\
\emph{Title}: Shape Learning from Euclidean Shadows: A Computational Geometric Framework\\
\emph{Role}: PI

\item NSF: Mathematical \& Physical Sciences (MPS), \textbf{Declined} \\
\emph{Title}: Geometric Approach to the Shape Reconstruction from Noisy Data\\
\emph{Role}: co-PI \\
\emph{PI}: Rafal Komendarczyk, Tulane University, New Orleans.
\fi

\item \textbf{Travel Grant}, University of California, Berkeley
\item \textbf{UGC-CSIR NET Research Fellowship}, India, June 2012
  
\item \textbf{TIFR Junior Research Fellowship} for pursuing Integrated PhD
studies at TIFR-CAM, Bangalore, India. August 2009.
  
\item \textbf{Secured grade ``A'' in SCIENCE TALENT SEARCH EXAMINATION}\\
  conducted by JATIYA VIJNAN PARISAD and INDIAN SCIENCE
  CONGRESS ASSOCIATION
\end{itemize}


%\let\thefootnote\relax\footnote{Last updated: \today}

\iflong


\fi


\section{Academic Services}
\begin{itemize}[leftmargin=*]
\item I organized two special sessions on ``Topological Data Analysis for Non-linear Dynamics'' and ``Topological and Geometric Shape Reconstruction'' at the Joint Mathematics Meetings (JMM) in Washington DC, 2025. 

\item I co-organized a special session on ``Advances in Applied Topology and Topological Data Analysis'' at the AMS southeastern sectional meeting at Tulane University in 2025. 

\item I have been a reviewer for journals, like Discrete and Computational Geometry, Foundations of Data Science, Journal of Combinatorial Optimization, and International Journal of Computational Geometry.
\item I have been a reviewer for conferences, like the International Symposium on
Computational Geometry, International Symposium on Spatial and Temporal
Databases, European Workshop on Computational Geometry, ACM International
Conference on Advances in Geographic Information Systems, WADS Algorithms and
Data Structures Symposium.
\item I organized SIAM Graduate Student Chapters at Tulane University. 
\item I have been organizing a data science webinar series at the University of
California, Berkeley. 
\end{itemize}

\iflong
\section{Computational Skills}
Java, C, R, Python, Ruby, JavaScript, SQL, Bash. \\
\fi

\section{Software Projects}
\begin{itemize}[leftmargin=*]
  \item \textbf{Simplicial Complexes in JS}
  \hfill\href{https://github.com/tdajs/simplicial-complex}{\footnotesize \tt
    GitHub}\\ 
JavaScript implementation of some of the widely used computations on simplicial
complexes. The library also implements the Smith Normal Form in order to compute
the homology groups of an abstract complex.
  
\item \textbf{Shape Reconstruction Visualization}
\hfill \href{https://smajhi.com/shape-reconstruction}{\footnotesize \tt WebApp}
| \href{https://github.com/sushovan4/shape-reconstruction}{\footnotesize \tt
    GitHub}\\ To complement my PhD research, I implemented my topological
  reconstruction algorithm for planar metric graphs in this library.  The
  library is written in JavaScript and made available to users as a
  web-app.\\ \emph{Skills}: JavaScript, HTML, CSS.
\end{itemize}

\section{Popular Media}
\begin{itemize}[leftmargin=*]
\item Interviewed by GWU Hatchet for an  
  \href{https://gwhatchet.com/2024/08/19/new-kids-on-the-block-first-year-professors-dish-on-joining-the-ranks-at-gw/}{\tt Article}
    \hfill \textbf{\footnotesize August 2024}
\end{itemize}


\section{Entrepreneurial Experience}
\begin{itemize}[leftmargin=*]
\item \textbf{Scimetric Edulabs Private Limited}
  \hfill \textbf{\footnotesize December 2012--April 2017}\\ Bangalore,
  India\\ \emph{Role}: \textbf{co-founder} and \textbf{director}\\ In this
  start-up venture, our objective was to motivate and train students in higher
  education. We won the franchise to work with several private colleges in India. We
  coached science students for standardized entrance tests for PhD and academic
  jobs. The company employed 6 trainers.
\end{itemize}

\iflong

\fi

\section{Hobby}
In my spare time, I write tutorials on \emph{random} topics in order to make
mathematics and statistics a little more interactive; they can be found here:
\url{https://www.smajhi.com/tutorials}. I also enjoy playing the piano and
classical guitar.

%\newpage
%\bibliographystyle{plain}

\end{document}

\section{Research Experience}
\begin{itemize}[leftmargin=*]
\item \textbf{Hausdorff vs Gromov--Hausdorff Distances on Manifolds} 
\hfill\textbf{\footnotesize April 2023--\emph{current}}\\ 
\emph{Collaborators}: \href{https://www.math.colostate.edu/~adams/}{Henry Adams}, University of Florida;
\href{https://www.math.cmu.edu/~ffrick/}{Florian Frick}, Carnegie Mellon University\\
In shape reconstruction from noisy samples, both the Hausdorff and
Gromov--Hausdorff distances are commonly used to facilitate noise models for the
sample. But the relationship between the two distance measures is not
well-understood. In this project, we further our understanding of their relationship, 
through the lenses of C\v{e}ch and Vietoris--Rips complexes.

\item \textbf{Vietoris--Rips Complexes near a Euclidean polytope} 
\hfill\textbf{\footnotesize June 2023--\emph{current}}\\ 
\emph{Collaborator}: \href{http://www.math.tulane.edu/~rako/}{Rafal Komendarczyk}, Tulane University\\
Finite reconstruction of Euclidean shapes that are not manifolds is a
challenging computational problem. The project is motivated by the problem of
topological reconstruction of Euclidean shapes that the stratified manifolds, 
examples include graphs, polytopes, etc. 

\item \textbf{Topological and Statistical Methods in Predicting the Crash\\ and
Recovery of Stock Markets} \hfill\textbf{\footnotesize March
2021--\emph{current}}\\ 
\emph{Collaborator}: \href{https://physics.nitsikkim.ac.in/profile.php?sn=23}{Md. Nurujjaman}, NIT Sikkim, India\\
In the aftermath of stock market crash due to COVID-19, not all sectors
recovered in the same way. We proposed novel models to capture the different
types of recovery profiles for Indian stocks. We also employed the
\textbf{Empirical Mode Decomposition} (EMD) for a statistical significance
analysis of our model. 

We currently look into the possibility of predicting a future crash in a
financial sector---using tools from \textbf{Topological Data Analysis} (TDA).

\item \textbf{Topological and Geometric Signature-Based Shape Comparison}
\hfill\textbf{\footnotesize October 2021--\emph{current}}\\ 
\emph{Collaborators}:
\href{https://cs.slu.edu/~chambers/}{Erin Chambers}, 
\href{https://elizabethmunch.com/}{Liz Munch}, and
\href{http://www.cs.tulane.edu/~carola/}{Carola Wenk}\\
We consider geometric and topological signatures to concisely represent large
datasets to facilitate their easy description and efficient comparison. To this
end, we look for signatures in the class of algebraic, geometric, and graphical
signatures. We have proposed new similarity measures for geometric graphs, and
studied their computational aspects.

\item \textbf{Computation of Gromov-Hausdorff Distance in Euclidean Space}
  \hfill\textbf{\footnotesize April
    2019--December 2020}\\ \emph{Collaborators}:
  \href{https://www.mi.fu-berlin.de/inf/groups/ag-ti/members/professoren/Alt_Helmut.html}{Helmut Alt},
   \href{https://www.cs.olemiss.edu/vitter/}{Jeffrey Vitter}, and
  \href{http://www.cs.tulane.edu/~carola/}{Carola Wenk}
\\ We investigate the computational aspects of the Gromov--Hausdorff distance between
sets equipped with the Euclidean metric. We used the Hausdorff distance under
isometry to develop an approximation algorithm for the Gromov--Hausdorff distance on
the real line with a tight approximation factor of $\left(1+\frac{1}{4}\right)$.

\item \textbf{Topological Reconstruction of Geodesic Spaces}
  \hfill \textbf{\footnotesize December 2016--May
    2019}\\ \emph{Collaborators}: 
  \href{https://www.cs.montana.edu/brittany/}{Brittany Fasy} and
  \href{http://129.81.170.14/~rako/}{Rafal Komendarczyk}\\
  \emph{Role}:
  Research Assistant\\
  \emph{PI}: \href{http://www.cs.tulane.edu/~carola/}{Carola Wenk} (NSF CCF-161846)\\
We investigate the reconstruction of geodesic subspaces of Euclidean spaces
using the Vietoris-Rips and \v{C}ech complexes from a dense sample around it. We
propose two new sampling parameters: \textbf{distortion} of embedding and
\textbf{convexity radius} of the underlying geodesic space. We guarantee a
successful computation of the Betti numbers. For the special case of planar
graphs, we also develop an algorithm for its geometric reconstruction.
  
\item \textbf{Dynamics and Prognosis of Chronic Myelogenous\\ Leukemia (CML)}
  \hfill \textbf{\footnotesize August 2012--November 2013}\\ National Center for
  Biological Sciences, TIFR, Bangalore, India\\ \emph{Role}: Junior Research
  Fellow\\ \emph{PI}: \href{https://www.icts.res.in/people/seema-nanda}{Seema
    Nanda}\\ In this joint effort to develop better prognostic tools for
  doctors, computational scientists teamed up with medical officers and
  biologists to understand the dynamics of CML by modeling the disease by
  systems of \textbf{differential equations}. In our parameter fitting, we made
  use of the big existing data collected from a large pool of CML patients. We
  also performed (statistical) \textbf{sensitivity analysis} to better
  understand the parameter spaces for our model.
\end{itemize}
